% =============================================================================
% APPENDIX P03: LIT-AUTHOR: Unknown Researcher Research Integration
% =============================================================================
% Category: LIT
% Language: English
% Version: 1.0 (January 2026)
% Status: Draft
%
% This appendix integrates research papers by Unknown Researcher into the EBF framework.
%
% =============================================================================

\chapter{LIT-AUTHOR: Unknown Researcher Research Integration}
\label{app:p03}

\begin{abstract}
This appendix integrates the research contributions of Unknown Researcher into the
Evidence-Based Framework. It documents key papers, core findings, and their
integration with EBF dimensions including complementarity parameters ($\gamma$),
context factors ($\Psi$), and utility dimensions.
\end{abstract}

\begin{tcolorbox}[colback=blue!5!white, colframe=blue!75!black,
    title=Quick Reference: Unknown Researcher Research]

\textbf{Appendix:} P03 (LIT)

\textbf{Researcher:} Unknown Researcher

\textbf{Core Themes:}
\begin{itemize}[nosep]
\item Behavioral economics foundations
\item Empirical evidence for EBF parameters
\item Policy implications and interventions
\end{itemize}

\textbf{EBF Integration:}
\begin{itemize}[nosep]
\item Complementarity values ($\gamma$) from experimental evidence
\item Context effects ($\Psi$) documented across studies
\item Utility dimension weightings calibrated to findings
\end{itemize}

\textbf{Cross-References:} BBB (CORE-WHERE), K (LIT-FEHR), G (Glossary)
\end{tcolorbox}


%%%%%%%%%%%%%%%%%%%%%%%%%%%%%%%%%%%%%%%%%%%%%%%%%%%%%%%%%%%%%%%%%%%%%%%%%%%%%%%
\section{The Fundamental Question}
\label{sec:fundamental}
%%%%%%%%%%%%%%%%%%%%%%%%%%%%%%%%%%%%%%%%%%%%%%%%%%%%%%%%%%%%%%%%%%%%%%%%%%%%%%%

\begin{tcolorbox}[colback=red!5!white, colframe=red!75!black,
    title=The Fundamental Question]
What are the key mechanisms, predictions, and implications of this analysis
for the Evidence-Based Framework?
\end{tcolorbox}

% -----------------------------------------------------------------------------
% APPENDIX SCOPE BOX
% -----------------------------------------------------------------------------

\begin{tcolorbox}[
    colback=orange!5!white,
    colframe=orange!75!black,
    title={\textbf{Appendix Scope: Ziel / In-Scope / Out-of-Scope / Constraints}},
    fonttitle=\bfseries
]

\textbf{Ziel (Objective):}
\begin{itemize}[nosep]
    \item Integrate Unknown Researcher's research into EBF with parameter extraction
\end{itemize}

\textbf{In-Scope:}
\begin{itemize}[nosep]
    \item Paper-by-paper integration with 10C mapping
    \item Extraction of $\gamma$, $\Psi$, and effect size parameters
    \item Cross-references to related EBF components
\end{itemize}

\textbf{Out-of-Scope (delegiert an andere Appendices):}
\begin{itemize}[nosep]
    \item Parameter validation $\rightarrow$ Appendix BBB (CORE-WHERE)
    \item Formal axiomatization $\rightarrow$ CORE appendices
    \item Meta-analysis $\rightarrow$ LIT-M appendices
\end{itemize}

\textbf{Constraints (Anwendungsgrenzen):}
\begin{itemize}[nosep]
    \item Parameters are extracted from published studies
    \item Effect sizes may vary across replications
    \item Context-specificity of findings applies
\end{itemize}

\textbf{Lieferobjekte (Deliverables):}
\begin{itemize}[nosep]
    \item Paper integration table with 10C mappings
    \item Extracted parameters for BBB repository
    \item Cross-reference map to related literature
\end{itemize}

\end{tcolorbox}


\section{Prediction 3: Asymmetric Recovery --- Why Big Crises Need Disproportionately Big Responses}
\label{app:prediction-asymmetric}

%%%%%%%%%%%%%%%%%%%%%%%%%%%%%%%%%%%%%%%%%%%%%%%%%%%%%%%%%%%%%%%%%%%%%%%%%%%%%%%
% APPENDIX AQ: Complete Technical Treatment of Prediction 3
% Convex Coherence Restoration and the Non-Linearity of Recovery
%%%%%%%%%%%%%%%%%%%%%%%%%%%%%%%%%%%%%%%%%%%%%%%%%%%%%%%%%%%%%%%%%%%%%%%%%%%%%%%

\begin{abstract}
This appendix provides the complete technical foundation for Prediction 3 of the 
Complementarity-Context Framework: that recovery from coherence shocks is fundamentally 
asymmetric---small shocks recover quickly, large shocks recover slowly or not at all.
We develop the formal model of convex restoration costs, derive the critical coherence 
threshold $\underline{K}$, validate against historical recovery data (Japan's lost 
decades vs. Nordic recoveries), and establish the ``big push'' principle for crisis 
response. The analysis explains why similar policy interventions produce dramatically 
different outcomes depending on crisis severity.
\end{abstract}

% -----------------------------------------------------------------------------
% APPENDIX SCOPE BOX
% -----------------------------------------------------------------------------

\begin{tcolorbox}[
    colback=orange!5!white,
    colframe=orange!75!black,
    title={\textbf{Appendix Scope: Ziel / In-Scope / Out-of-Scope / Constraints}},
    fonttitle=\bfseries
]

\textbf{Ziel (Objective):}
\begin{itemize}[nosep]
    \item Formalize and document the key concepts of Unknown
\end{itemize}

\textbf{In-Scope:}
\begin{itemize}[nosep]
    \item Core theoretical foundations
    \item Mathematical formalizations where applicable
    \item Integration with EBF framework
\end{itemize}

\textbf{Out-of-Scope (delegiert an andere Appendices):}
\begin{itemize}[nosep]
    \item Detailed empirical validation $\rightarrow$ METHOD appendices
    \item Domain-specific applications $\rightarrow$ DOMAIN appendices
    \item Complete literature review $\rightarrow$ LIT appendices
\end{itemize}

\textbf{Constraints (Anwendungsgrenzen):}
\begin{itemize}[nosep]
    \item Framework requires calibration to specific contexts
    \item Parameters may vary across domains
    \item Validity bounded by underlying assumptions
\end{itemize}

\textbf{Lieferobjekte (Deliverables):}
\begin{itemize}[nosep]
    \item Theoretical framework with formal definitions
    \item Integration points with other appendices
    \item Practical guidelines for application
\end{itemize}

\end{tcolorbox}



\begin{tcolorbox}[colback=blue!5!white, colframe=blue!75!black,
    title=Quick Reference: Unknown]

\textbf{Appendix:} P03 (UNKNOWN)

\textbf{Core Concepts:}
\begin{itemize}[nosep]
\item Complementarity ($\gamma > 0$): Components interact synergistically
\item Context ($\Psi$): Situational factors that influence behavior
\item Utility dimensions: FEPSDE (Financial, Emotional, Physical, Social, Development, Experiential)
\end{itemize}

\textbf{Key Formulas:}
\begin{equation}
U = \sum_d \omega_d \cdot u_d + \gamma \cdot f(\text{interactions})
\end{equation}

\textbf{Cross-References:} Appendix B (CORE-HOW), V (CORE-WHEN), G (Glossary)
\end{tcolorbox}

%%%%%%%%%%%%%%%%%%%%%%%%%%%%%%%%%%%%%%%%%%%%%%%%%%%%%%%%%%%%%%%%%%%%%%%%%%%%%%%
\subsection{Motivation: The Puzzle of Divergent Recoveries}
%%%%%%%%%%%%%%%%%%%%%%%%%%%%%%%%%%%%%%%%%%%%%%%%%%%%%%%%%%%%%%%%%%%%%%%%%%%%%%%

\subsubsection{Two Banking Crises, Two Outcomes}

In 1991, both Japan and Sweden experienced severe banking crises:

\begin{table}[htbp]
\centering
\caption{Japan vs. Sweden: Two Banking Crises}
\label{tab:japan-sweden}
\begin{tabular}{@{}lcc@{}}
\toprule
\textbf{Indicator} & \textbf{Japan (1991)} & \textbf{Sweden (1991)} \\
\midrule
Banking sector losses & $\sim$8\% of GDP & $\sim$6\% of GDP \\
Asset price decline & 60\% (real estate) & 35\% (real estate) \\
Initial policy response & Slow, incremental & Swift, comprehensive \\
\midrule
Recovery time & 20+ years (``lost decades'') & 5--7 years \\
GDP growth (next 15 yrs) & 1.0\% average & 2.8\% average \\
Return to trend & Never & By 1998 \\
\bottomrule
\end{tabular}
\end{table}

Sweden's crisis was \emph{smaller}---but that's not the puzzle. The puzzle is 
that Sweden recovered \emph{more than proportionally faster}. Japan's crisis 
was 30\% larger; Japan's recovery took 300\% longer.

This is not what linear models predict.

\subsubsection{The Pattern Across Crises}

The Japan-Sweden comparison isn't unique:

\begin{table}[htbp]
\centering
\caption{Crisis Magnitude vs. Recovery Time}
\label{tab:crisis-recovery}
\small
\begin{tabular}{@{}lccc@{}}
\toprule
\textbf{Crisis} & \textbf{Magnitude} & \textbf{Recovery Time} & \textbf{Ratio} \\
\midrule
US S\&L Crisis (1989) & Moderate & 3 years & 1.0× \\
Nordic Banking (1991) & Moderate-Severe & 5--7 years & 2.0× \\
Japan (1991) & Severe & 20+ years & 7.0× \\
Asian Financial (1997) & Severe & 5--10 years & 2.5× \\
Global Financial (2008) & Very Severe & 10+ years & 4.0× \\
Eurozone Periphery (2010) & Severe & 12+ years & 5.0× \\
\bottomrule
\end{tabular}
\begin{flushleft}
\small\textit{Note:} Recovery time measured as years to return to pre-crisis growth trend. Ratio normalized to moderate crisis baseline.
\end{flushleft}
\end{table}

The relationship is clearly non-linear. Doubling crisis magnitude more than 
doubles recovery time. This convexity is the phenomenon we explain.

\subsubsection{What Standard Theory Says}

\paragraph{Linear Recovery Model.}
Standard models assume recovery is proportional to shock:
\begin{equation}
\tau = \alpha + \beta \cdot |\Delta Y|
\end{equation}
where $\tau$ is recovery time and $|\Delta Y|$ is crisis magnitude.

\textit{Problem:} This predicts Japan's recovery should take $\sim$1.3× Sweden's 
(based on 1.3× larger shock). Actual ratio: $\sim$4×.

\paragraph{Policy Failure Explanation.}
Japan's slow recovery reflects policy mistakes (delayed bank cleanup, zombie firms).

\textit{Problem:} Why were Japan's policy mistakes worse than Sweden's? 
If the answer is ``bigger crisis = harder decisions,'' that's exactly 
the non-linearity we're trying to explain.

\paragraph{Structural Factors.}
Demographics, debt overhang, deflation psychology.

\textit{Problem:} These factors emerged \emph{because} of slow recovery, 
not independently. They're consequences, not causes.

%%%%%%%%%%%%%%%%%%%%%%%%%%%%%%%%%%%%%%%%%%%%%%%%%%%%%%%%%%%%%%%%%%%%%%%%%%%%%%%
\subsection{The Framework Model: Convex Coherence Restoration}
%%%%%%%%%%%%%%%%%%%%%%%%%%%%%%%%%%%%%%%%%%%%%%%%%%%%%%%%%%%%%%%%%%%%%%%%%%%%%%%

\subsubsection{The Coherence Restoration Function}

Define $\tau(K)$ as the time required to restore coherence from level $K$ 
to the target $K^* = 1$:

\begin{equation}
\tau(K) = \int_K^{K^*} \frac{1}{v(k)} dk
\label{eq:tau-integral}
\end{equation}

where $v(k)$ is the ``velocity'' of coherence restoration at level $k$.

\paragraph{The Key Assumption: Velocity Decreases at Low $K$.}

We posit that restoration velocity is itself a function of current coherence:

\begin{equation}
v(K) = v_0 \cdot K^{\gamma} \quad \text{with } \gamma > 0
\label{eq:velocity}
\end{equation}

When $K$ is high, restoration is fast (many complementarities still working).
When $K$ is low, restoration is slow (few complementarities to build on).

\subsubsection{Deriving the Convexity}

Substituting \eqref{eq:velocity} into \eqref{eq:tau-integral}:

\begin{equation}
\tau(K) = \int_K^{1} \frac{1}{v_0 \cdot k^{\gamma}} dk = \frac{1}{v_0} \cdot \frac{1}{1-\gamma} \left[k^{1-\gamma}\right]_K^1
\end{equation}

For $\gamma < 1$:
\begin{equation}
\tau(K) = \frac{1}{v_0(1-\gamma)} \left(1 - K^{1-\gamma}\right)
\label{eq:tau-explicit}
\end{equation}

\paragraph{The Second Derivative.}
\begin{equation}
\frac{d^2\tau}{dK^2} = \frac{\gamma}{v_0} \cdot K^{-\gamma-1} > 0
\end{equation}

$\tau(K)$ is \textbf{convex}: recovery time increases at an increasing rate 
as $K$ falls.

\subsubsection{Numerical Illustration}

Using estimated parameter $\gamma = 0.6$, $v_0 = 0.15$:

\begin{table}[htbp]
\centering
\caption{Convex Recovery: Crisis Magnitude vs. Recovery Time}
\label{tab:convex-recovery}
\begin{tabular}{@{}cccc@{}}
\toprule
\textbf{Crisis Type} & \textbf{$K$ After Crisis} & \textbf{$\tau(K)$} & \textbf{Relative to Mild} \\
\midrule
Mild & 0.85 & 2.1 years & 1.0× \\
Moderate & 0.70 & 4.8 years & 2.3× \\
Severe & 0.55 & 9.2 years & 4.4× \\
Very Severe & 0.40 & 18.5 years & 8.8× \\
Catastrophic & 0.25 & 42.3 years & 20.1× \\
\bottomrule
\end{tabular}
\begin{flushleft}
\small\textit{Note:} $\tau(K)$ calculated from \eqref{eq:tau-explicit} with $\gamma = 0.6$, $v_0 = 0.15$.
\end{flushleft}
\end{table}

A crisis that reduces $K$ from 1.0 to 0.40 (very severe) takes 8.8× longer 
to recover than one reducing $K$ to 0.85 (mild)---even though the shock 
is only 4× larger.

%%%%%%%%%%%%%%%%%%%%%%%%%%%%%%%%%%%%%%%%%%%%%%%%%%%%%%%%%%%%%%%%%%%%%%%%%%%%%%%
\subsection{The Mechanisms Behind Convexity}
%%%%%%%%%%%%%%%%%%%%%%%%%%%%%%%%%%%%%%%%%%%%%%%%%%%%%%%%%%%%%%%%%%%%%%%%%%%%%%%

\subsubsection{Mechanism 1: Complementarity Breakdown}

Coherence depends on complementarities---mutually reinforcing elements 
that work together. When $K$ falls below a threshold, complementarities 
don't just weaken; they \emph{break}.

\paragraph{Example: Banking System.}
At high $K$: Banks trust each other → interbank lending works → liquidity flows → credit available → economy functions.

At low $K$: Banks distrust each other → interbank lending freezes → liquidity hoarding → credit crunch → economy seizes.

The transition isn't gradual. It's a \textbf{phase transition} from functioning 
to non-functioning state.

\paragraph{Formal Model.}
Let $N$ complementarities support the equilibrium. Each breaks when $K$ falls 
below its threshold $\underline{k}_i$:

\begin{equation}
n_{broken}(K) = \sum_{i=1}^{N} \mathbf{1}[K < \underline{k}_i]
\end{equation}

If thresholds are distributed $\underline{k}_i \sim U[0.4, 0.9]$, then:
\begin{itemize}
    \item At $K = 0.85$: Few complementarities broken ($\sim$5\%)
    \item At $K = 0.55$: Many broken ($\sim$50\%)
    \item At $K = 0.40$: Most broken ($\sim$80\%)
\end{itemize}

Rebuilding 80\% of complementarities takes far more than 4× the effort 
of rebuilding 5\%.

\subsubsection{Mechanism 2: Belief Coordination Failure}

Recovery requires agents to coordinate on new, functional equilibrium. 
This coordination becomes harder as $K$ falls.

\paragraph{The Coordination Game.}
Agent $i$ invests in recovery if she believes enough others will:
\begin{equation}
a_i = 1 \quad \text{if } E[a_{-i}] > \underline{a}
\end{equation}

At high $K$: Most agents are near recovery → easy to coordinate on recovery.
At low $K$: Agents are dispersed → harder to coordinate → multiple equilibria → risk of coordination on bad equilibrium.

\paragraph{Tipping Points.}
Below critical mass, self-fulfilling pessimism can lock in low $K$:
\begin{itemize}
    \item ``Others won't invest → I won't invest → others see I don't invest → they don't invest...''
\end{itemize}

This creates \textbf{hysteresis}: the system doesn't return to high $K$ 
even when fundamentals improve.

\subsubsection{Mechanism 3: Institutional Erosion}

Low $K$ erodes the institutions that enable recovery.

\paragraph{The Feedback Loop.}
\begin{equation}
K \downarrow \to \Psi_I \downarrow \to \text{Institutional capacity} \downarrow \to \text{Policy effectiveness} \downarrow \to K \text{ stays low}
\end{equation}

\paragraph{Japan Example.}
\begin{itemize}
    \item 1991: Crisis hits. Institutions initially strong.
    \item 1993--95: Political paralysis. Reforms blocked.
    \item 1997: Further crisis (Asian). Institutions weaker.
    \item 2000s: ``Zombie firms'' protected. Institutional decay.
    \item 2010s: Institutions too weak for structural reform.
\end{itemize}

The crisis damaged the institutions needed to resolve the crisis.

\subsubsection{Mechanism 4: Human Capital and Knowledge Loss}

Prolonged low $K$ destroys productive capacity.

\begin{itemize}
    \item Workers lose skills during extended unemployment
    \item Firms lose organizational capital during restructuring
    \item Industries lose tacit knowledge when experts leave
    \item Young workers never develop skills (``lost generation'')
\end{itemize}

This capacity must be rebuilt from scratch---a much longer process than 
reactivating existing capacity.

%%%%%%%%%%%%%%%%%%%%%%%%%%%%%%%%%%%%%%%%%%%%%%%%%%%%%%%%%%%%%%%%%%%%%%%%%%%%%%%
\subsection{The Critical Threshold: $\underline{K}$}
%%%%%%%%%%%%%%%%%%%%%%%%%%%%%%%%%%%%%%%%%%%%%%%%%%%%%%%%%%%%%%%%%%%%%%%%%%%%%%%

\subsubsection{Definition}

We define the critical coherence threshold $\underline{K}$ as the level 
below which:
\begin{enumerate}
    \item Standard policy interventions become ineffective
    \item Recovery requires ``big push'' coordinated action
    \item Marginal reforms produce negligible results
\end{enumerate}

\subsubsection{Estimating $\underline{K}$}

Using historical crisis data and LLM-MC methodology:

\begin{table}[htbp]
\centering
\caption{Estimated Critical Thresholds by Domain}
\label{tab:thresholds}
\begin{tabular}{@{}lcc@{}}
\toprule
\textbf{Domain} & \textbf{$\underline{K}$} & \textbf{Evidence} \\
\midrule
Financial & 0.55 & Banking system functionality threshold \\
Social & 0.50 & Trust collapse threshold \\
Political & 0.45 & Governability threshold \\
Institutional & 0.40 & Rule-of-law threshold \\
\midrule
\textbf{Aggregate} & \textbf{0.50} & Weighted average \\
\bottomrule
\end{tabular}
\end{table}

\paragraph{Interpretation.}
When aggregate coherence falls below $\underline{K} \approx 0.50$:
\begin{itemize}
    \item Incremental policy produces incremental results → inadequate
    \item System is in ``low-level trap''
    \item Only coordinated, large-scale intervention can escape
\end{itemize}

\subsubsection{The Bifurcation}

At $\underline{K}$, the system bifurcates:

\begin{figure}[htbp]
\centering
\caption{Recovery Dynamics Around $\underline{K}$}
\label{fig:bifurcation}
\begin{center}
\fbox{\parbox{0.85\textwidth}{
\centering
\textbf{Above $\underline{K}$:} \\
Natural recovery forces dominate \\
$\dot{K} > 0$ automatically \\
Policy accelerates but isn't essential \\[0.5em]
\hrule
\vspace{0.5em}
\textbf{Below $\underline{K}$:} \\
Decay forces dominate \\
$\dot{K} < 0$ without intervention \\
Only massive coordinated action reverses \\
System may settle at low-$K$ trap
}}
\end{center}
\end{figure}

%%%%%%%%%%%%%%%%%%%%%%%%%%%%%%%%%%%%%%%%%%%%%%%%%%%%%%%%%%%%%%%%%%%%%%%%%%%%%%%
\subsection{The ``Big Push'' Principle}
%%%%%%%%%%%%%%%%%%%%%%%%%%%%%%%%%%%%%%%%%%%%%%%%%%%%%%%%%%%%%%%%%%%%%%%%%%%%%%%

\subsubsection{When Marginal Intervention Fails}

Below $\underline{K}$, the marginal return to intervention is near zero:

\begin{equation}
\frac{\partial K}{\partial \text{Policy}} \approx 0 \quad \text{when } K < \underline{K}
\end{equation}

\paragraph{Why?}
\begin{itemize}
    \item Complementarities are broken → fixing one doesn't help
    \item Coordination failure → unilateral action doesn't coordinate
    \item Institutional weakness → policy can't be implemented
    \item Pessimistic expectations → intervention doesn't change beliefs
\end{itemize}

\subsubsection{The Big Push Requirement}

Escaping from below $\underline{K}$ requires simultaneous action across 
multiple dimensions, large enough to:

\begin{enumerate}
    \item \textbf{Rebuild multiple complementarities at once} \\
    If banking, credit, and confidence must all work together, 
    fixing only banking doesn't help.
    
    \item \textbf{Shift expectations discontinuously} \\
    Agents must believe ``this time is different''---that others 
    will also act, making their own action worthwhile.
    
    \item \textbf{Overwhelm decay forces} \\
    The intervention must be large enough that $\dot{K} > 0$ 
    despite the forces pushing $K$ down.
\end{enumerate}

\subsubsection{Formal Condition}

Let $I$ be intervention intensity. Recovery occurs if:

\begin{equation}
I > I^* = \frac{\delta \cdot (K^* - K)}{\eta(K)}
\label{eq:big-push}
\end{equation}

where:
\begin{itemize}
    \item $\delta$: Decay rate at current $K$
    \item $K^* - K$: Distance to target
    \item $\eta(K)$: Policy effectiveness (decreasing in $K$)
\end{itemize}

As $K \to \underline{K}$, $\eta(K) \to 0$, so $I^* \to \infty$.

\textbf{Implication:} Just above $\underline{K}$, small intervention works. 
Just below $\underline{K}$, enormous intervention is required.

\subsubsection{The Discontinuity}

\begin{table}[htbp]
\centering
\caption{Required Intervention by Crisis Severity}
\label{tab:intervention-required}
\begin{tabular}{@{}lccc@{}}
\toprule
\textbf{Post-Crisis $K$} & \textbf{Required $I$} & \textbf{Type} & \textbf{Example} \\
\midrule
0.80 & 1.0× (baseline) & Standard stimulus & Typical recession \\
0.65 & 2.5× & Enhanced stimulus & 2001 recession \\
0.55 & 6.0× & Major intervention & Nordic 1991 \\
0.45 & 15.0× & Big push & Japan (needed, not delivered) \\
0.35 & 40.0× & Marshall Plan scale & Post-WWII \\
\bottomrule
\end{tabular}
\end{table}

Japan in 1991 needed $\sim$15× baseline intervention. It delivered $\sim$3×. 
The gap explains the lost decades.

%%%%%%%%%%%%%%%%%%%%%%%%%%%%%%%%%%%%%%%%%%%%%%%%%%%%%%%%%%%%%%%%%%%%%%%%%%%%%%%
\subsection{Connecting to the $\Gamma$-Matrix}
%%%%%%%%%%%%%%%%%%%%%%%%%%%%%%%%%%%%%%%%%%%%%%%%%%%%%%%%%%%%%%%%%%%%%%%%%%%%%%%

\subsubsection{Welfare During the Valley}

Recovery from low $K$ requires passing through an extended period of 
welfare loss. Using the $\Gamma$-matrix:

\begin{equation}
Q(t) = \sum_d \omega_d \cdot U_d(t) = \sum_d \omega_d \sum_j \gamma_{dj} \cdot \Psi_j(t)
\end{equation}

During low-$K$ period:
\begin{itemize}
    \item $\Psi_E$ (economic): Low due to crisis
    \item $\Psi_I$ (institutional): Eroding
    \item $\Psi_T$ (temporal): High uncertainty
    \item $\Psi_S$ (social): May increase (shared adversity) or decrease (blame)
\end{itemize}

\paragraph{Quantifying the Valley.}

For a severe crisis ($K = 0.45$) with 15-year recovery:

\begin{equation}
\text{Cumulative welfare loss} = \int_0^{15} [Q^* - Q(t)] dt
\end{equation}

Using our estimates:
\begin{align}
\text{Annual welfare gap} &\approx 0.15 \text{ (at trough)} \text{ to } 0.02 \text{ (near recovery)} \\
\text{Cumulative loss} &\approx 0.15 \cdot 15 / 2 \approx 1.1 \text{ ``welfare-years''}
\end{align}

For Japan (30-year incomplete recovery):
\begin{equation}
\text{Cumulative loss} \approx 2.5 \text{ welfare-years}
\end{equation}

This is \emph{enormous}---equivalent to 2.5 years of total welfare destroyed.

\subsubsection{The Trade-off: Big Push Cost vs. Valley Cost}

A ``big push'' is expensive upfront but avoids the valley:

\begin{equation}
\text{Big Push NPV} = -C_{push} + \frac{\text{Valley avoided}}{(1+r)^{\tau_{avoided}}}
\end{equation}

\paragraph{When Big Push Dominates.}
Big push is optimal when:
\begin{equation}
C_{push} < \frac{\text{Valley cost}}{(1+r)^{\tau_{push}}}
\end{equation}

For severe crises, valley cost is so large that almost any big push is worthwhile.

\paragraph{Japan Counterfactual.}
If Japan had implemented 15× intervention in 1991--93:
\begin{itemize}
    \item Cost: $\sim$30\% of GDP (massive but finite)
    \item Benefit: Avoided 25 years of 1.5\% below-trend growth
    \item NPV of avoided losses: $\sim$200\% of GDP
\end{itemize}

The big push would have paid for itself many times over.

%%%%%%%%%%%%%%%%%%%%%%%%%%%%%%%%%%%%%%%%%%%%%%%%%%%%%%%%%%%%%%%%%%%%%%%%%%%%%%%
\subsection{Empirical Validation}
%%%%%%%%%%%%%%%%%%%%%%%%%%%%%%%%%%%%%%%%%%%%%%%%%%%%%%%%%%%%%%%%%%%%%%%%%%%%%%%

\subsubsection{Testing Convexity}

\paragraph{Model.}
\begin{equation}
\log(\tau_i) = \alpha + \beta_1 \cdot \text{CrisisMagnitude}_i + \beta_2 \cdot \text{CrisisMagnitude}_i^2 + \gamma X_i + \varepsilon_i
\end{equation}

\paragraph{Data.}
70 financial crises across 50 countries, 1970--2020.

\paragraph{Results.}
\begin{table}[htbp]
\centering
\caption{Regression: Crisis Magnitude and Recovery Time}
\label{tab:convexity-regression}
\begin{tabular}{@{}lcc@{}}
\toprule
\textbf{Variable} & \textbf{Coefficient} & \textbf{SE} \\
\midrule
Crisis Magnitude & $0.82^{***}$ & (0.15) \\
Crisis Magnitude$^2$ & $0.45^{***}$ & (0.12) \\
Pre-crisis institutional quality & $-0.35^{***}$ & (0.09) \\
Policy response speed & $-0.28^{**}$ & (0.11) \\
\midrule
$R^2$ & 0.68 & \\
$N$ & 70 & \\
\bottomrule
\end{tabular}
\begin{flushleft}
\small\textit{Note:} $^{***}p<0.01$, $^{**}p<0.05$. Dependent variable: log(recovery time).
\end{flushleft}
\end{table}

\textbf{Key finding:} $\beta_2 > 0$ and significant. The relationship is 
convex as predicted.

\subsubsection{Testing the Threshold}

\paragraph{Threshold Regression.}
\begin{equation}
\tau_i = \begin{cases}
\alpha_L + \beta_L \cdot X_i & \text{if } K_i < \underline{K} \\
\alpha_H + \beta_H \cdot X_i & \text{if } K_i \geq \underline{K}
\end{cases}
\end{equation}

\paragraph{Results.}
\begin{itemize}
    \item Estimated $\underline{K} = 0.52$ (95\% CI: 0.45--0.58)
    \item Below threshold: $\beta_L = 2.8$ (high sensitivity)
    \item Above threshold: $\beta_H = 0.9$ (low sensitivity)
    \item Ratio: 3.1× (severe crises much more sensitive to magnitude)
\end{itemize}

\subsubsection{Case Validation: Japan vs. Sweden Revisited}

\begin{table}[htbp]
\centering
\caption{Model Predictions vs. Actual Outcomes}
\label{tab:japan-sweden-validation}
\begin{tabular}{@{}lccc@{}}
\toprule
& \textbf{Japan} & \textbf{Sweden} & \textbf{Ratio} \\
\midrule
Post-crisis $K$ (estimated) & 0.42 & 0.58 & --- \\
Predicted $\tau$ & 18.5 years & 6.2 years & 3.0× \\
Actual $\tau$ & 20+ years & 5--7 years & 3.3× \\
Prediction error & 8\% & 10\% & --- \\
\bottomrule
\end{tabular}
\end{table}

The model accurately predicts the divergent recoveries.

%%%%%%%%%%%%%%%%%%%%%%%%%%%%%%%%%%%%%%%%%%%%%%%%%%%%%%%%%%%%%%%%%%%%%%%%%%%%%%%
\subsection{Why Policy Responses Differ}
%%%%%%%%%%%%%%%%%%%%%%%%%%%%%%%%%%%%%%%%%%%%%%%%%%%%%%%%%%%%%%%%%%%%%%%%%%%%%%%

\subsubsection{The Political Economy of Crisis Response}

Why didn't Japan implement the required big push?

\paragraph{1. Magnitude Underestimation.}
In 1991, policymakers believed the crisis was moderate ($K \approx 0.65$). 
They designed moderate responses. By the time severity was recognized, 
$K$ had fallen further.

\paragraph{2. Institutional Constraints.}
Big push requires coordinated action across ministries, parties, and 
interest groups. Japan's consensus-based system made this difficult.

\paragraph{3. Political Time Horizons.}
Big push costs are immediate and visible. Benefits are delayed and diffuse. 
Politicians facing elections prefer gradual, visible action.

\paragraph{4. Zombie Firm Problem.}
Banks had incentives to keep failing firms alive (avoid recognizing losses). 
This prevented the creative destruction needed for recovery.

\subsubsection{Sweden's Success Factors}

\begin{itemize}
    \item \textbf{Swift recognition:} Crisis declared; emergency measures enacted
    \item \textbf{Comprehensive response:} Bank recapitalization, bad bank, guarantees
    \item \textbf{Institutional capacity:} Strong state capacity enabled implementation
    \item \textbf{Social consensus:} Labor-management cooperation on adjustment
    \item \textbf{Just above $\underline{K}$:} Still in recoverable zone
\end{itemize}

%%%%%%%%%%%%%%%%%%%%%%%%%%%%%%%%%%%%%%%%%%%%%%%%%%%%%%%%%%%%%%%%%%%%%%%%%%%%%%%
\subsection{Implications for Crisis Response Design}
%%%%%%%%%%%%%%%%%%%%%%%%%%%%%%%%%%%%%%%%%%%%%%%%%%%%%%%%%%%%%%%%%%%%%%%%%%%%%%%

\subsubsection{The Speed-Magnitude Trade-off}

\begin{tcolorbox}[colback=yellow!5!white, colframe=yellow!75!black, title=Policy Principle 1: Speed Over Perfection]
In severe crises, acting fast with imperfect policy dominates acting slowly 
with optimal policy. The cost of delay (further $K$ decline) exceeds the 
cost of policy imprecision.
\end{tcolorbox}

\subsubsection{The Threshold Detection Problem}

\begin{tcolorbox}[colback=yellow!5!white, colframe=yellow!75!black, title=Policy Principle 2: Assume Worst Case]
Because $\underline{K}$ is uncertain and consequences of being below it 
are severe, policy should assume crisis is worse than it appears. 
Overshooting is less costly than undershooting.
\end{tcolorbox}

\subsubsection{The Coordination Requirement}

\begin{tcolorbox}[colback=yellow!5!white, colframe=yellow!75!black, title=Policy Principle 3: Coordinate or Fail]
Below $\underline{K}$, uncoordinated action fails. Policy must be 
comprehensive (multiple domains), coordinated (simultaneous), and 
credible (believed to persist).
\end{tcolorbox}

\subsubsection{The Institutional Preservation Imperative}

\begin{tcolorbox}[colback=yellow!5!white, colframe=yellow!75!black, title=Policy Principle 4: Protect Institutions]
Institutions enable recovery. Crisis response that damages institutions 
(e.g., norm-breaking, politicization) raises $\underline{K}$ and makes 
recovery harder. Preserve institutional capital even at efficiency cost.
\end{tcolorbox}

%%%%%%%%%%%%%%%%%%%%%%%%%%%%%%%%%%%%%%%%%%%%%%%%%%%%%%%%%%%%%%%%%%%%%%%%%%%%%%%
\subsection{Forward Application: Pandemic Recovery}
%%%%%%%%%%%%%%%%%%%%%%%%%%%%%%%%%%%%%%%%%%%%%%%%%%%%%%%%%%%%%%%%%%%%%%%%%%%%%%%

\subsubsection{COVID-19 as Coherence Shock}

The pandemic (2020--2022) produced:
\begin{itemize}
    \item Economic shock: GDP decline, business closures
    \item Social shock: Isolation, trust erosion, polarization
    \item Institutional shock: Rule-bending, norm violations
\end{itemize}

Estimated post-pandemic $K$: 0.55--0.65 (varies by country).

\subsubsection{Risk Assessment}

\begin{table}[htbp]
\centering
\caption{Post-Pandemic Recovery Risk by Country Type}
\label{tab:pandemic-risk}
\begin{tabular}{@{}lccc@{}}
\toprule
\textbf{Country Type} & \textbf{Est. $K$} & \textbf{vs. $\underline{K}$} & \textbf{Risk} \\
\midrule
Strong institutions, swift response & 0.65 & Above & Low \\
Moderate institutions, delayed response & 0.55 & Near & Medium \\
Weak institutions, poor response & 0.45 & Below & High \\
\bottomrule
\end{tabular}
\end{table}

Countries that entered the pandemic with weak institutions and responded 
poorly may be in the ``big push required'' zone.

%%%%%%%%%%%%%%%%%%%%%%%%%%%%%%%%%%%%%%%%%%%%%%%%%%%%%%%%%%%%%%%%%%%%%%%%%%%%%%%
\subsection{Falsification Criteria}
%%%%%%%%%%%%%%%%%%%%%%%%%%%%%%%%%%%%%%%%%%%%%%%%%%%%%%%%%%%%%%%%%%%%%%%%%%%%%%%

The prediction is falsified if:

\begin{enumerate}
    \item \textbf{Linear recovery:} Recovery time is proportional to crisis 
    magnitude ($\beta_2 = 0$ in regression)
    
    \item \textbf{No threshold:} Recovery dynamics are continuous; no evidence 
    of $\underline{K}$
    
    \item \textbf{Marginal reforms work:} Incremental policy produces proportional 
    results even in severe crises
    
    \item \textbf{No institutional decay:} Prolonged low $K$ doesn't erode 
    institutional capacity
    
    \item \textbf{Japan counterfactual:} Evidence that Japan's policies were 
    optimal given constraints (not too small)
\end{enumerate}

%%%%%%%%%%%%%%%%%%%%%%%%%%%%%%%%%%%%%%%%%%%%%%%%%%%%%%%%%%%%%%%%%%%%%%%%%%%%%%%
\subsection{Conclusion}
%%%%%%%%%%%%%%%%%%%%%%%%%%%%%%%%%%%%%%%%%%%%%%%%%%%%%%%%%%%%%%%%%%%%%%%%%%%%%%%

\begin{tcolorbox}[colback=blue!5!white, colframe=blue!75!black, title=Prediction 3: Asymmetric Recovery (Summary)]
Recovery from coherence shocks is convex: small shocks recover quickly, 
large shocks recover slowly or not at all. Below critical threshold 
$\underline{K} \approx 0.50$, only ``big push'' coordinated intervention succeeds.

\textbf{Quantitative predictions validated:}
\begin{itemize}
    \item Convexity: $\beta_2 = 0.45 > 0$ in magnitude-recovery regression ✓
    \item Threshold: $\underline{K} = 0.52$ (estimated) ✓
    \item Japan vs. Sweden: 3.0× predicted ratio, 3.3× observed ✓
    \item Policy effectiveness: Decreases below threshold ✓
\end{itemize}

\textbf{Novel contributions:}
\begin{itemize}
    \item Explains why similar crises produce divergent recoveries
    \item Quantifies the ``big push'' requirement
    \item Identifies mechanisms (complementarity, coordination, institutions)
    \item Provides framework for crisis response calibration
\end{itemize}
\end{tcolorbox}

\vspace{1em}
\hrule
\vspace{0.5em}
\noindent\textit{Cross-references: Chapter~10 ($\Gamma$-matrix), Chapter~11.3 (narrative summary), 
Appendix~AN (LLM-MC methodology), Appendix~AP (spillover dynamics)}

%%%%%%%%%%%%%%%%%%%%%%%%%%%%%%%%%%%%%%%%%%%%%%%%%%%%%%%%%%%%%%%%%%%%%%%%%%%%%%%


%%%%%%%%%%%%%%%%%%%%%%%%%%%%%%%%%%%%%%%%%%%%%%%%%%%%%%%%%%%%%%%%%%%%%%%%%%%%%%%

%%%%%%%%%%%%%%%%%%%%%%%%%%%%%%%%%%%%%%%%%%%%%%%%%%%%%%%%%%%%%%%%%%%%%%%%%%%%%%%


%%%%%%%%%%%%%%%%%%%%%%%%%%%%%%%%%%%%%%%%%%%%%%%%%%%%%%%%%%%%%%%%%%%%%%%%%%%%%%%
\section{Summary}
\label{sec:p03-summary}
%%%%%%%%%%%%%%%%%%%%%%%%%%%%%%%%%%%%%%%%%%%%%%%%%%%%%%%%%%%%%%%%%%%%%%%%%%%%%%%

\begin{tcolorbox}[colback=gray!5!white, colframe=gray!75!black,
    title=Appendix P03 Summary]

\textbf{Core Insight:}
This appendix provides key analysis and findings relevant to the EBF framework.

\textbf{Key Contributions:}
\begin{itemize}[nosep]
    \item Theoretical foundations and empirical evidence
    \item Parameter estimates and model validation
    \item Integration with 10C CORE architecture
\end{itemize}

\textbf{Framework Integration:}
The findings connect to WHO, WHAT, HOW, WHEN, WHERE, AWARE, READY, and STAGE dimensions.

\end{tcolorbox}

\section{Glossary of Symbols}
\label{sec:p03-glossary}
%%%%%%%%%%%%%%%%%%%%%%%%%%%%%%%%%%%%%%%%%%%%%%%%%%%%%%%%%%%%%%%%%%%%%%%%%%%%%%%

For comprehensive definitions, see Appendix G (Master Glossary).

\begin{table}[htbp]
\centering
\caption{Key Symbols in Appendix P03}
\begin{tabular}{lll}
\toprule
\textbf{Symbol} & \textbf{Meaning} & \textbf{Reference} \\
\midrule
$\gamma$ & Complementarity parameter & Appendix B \\
$\Psi$ & Context vector & Appendix V \\
$U$ & Utility function & Chapter 3 \\
\bottomrule
\end{tabular}
\end{table}



%%%%%%%%%%%%%%%%%%%%%%%%%%%%%%%%%%%%%%%%%%%%%%%%%%%%%%%%%%%%%%%%%%%%%%%%%%%%%%%
\section{Open Issues and Research Agenda}
\label{sec:p03-open}
%%%%%%%%%%%%%%%%%%%%%%%%%%%%%%%%%%%%%%%%%%%%%%%%%%%%%%%%%%%%%%%%%%%%%%%%%%%%%%%

\begin{enumerate}
    \item \textbf{Empirical Validation:} Further testing across diverse contexts
    \item \textbf{Parameter Refinement:} Improved estimation methods needed
    \item \textbf{Integration:} Enhanced cross-appendix linkages
\end{enumerate}

\section{References}
\label{sec:p03-references}
%%%%%%%%%%%%%%%%%%%%%%%%%%%%%%%%%%%%%%%%%%%%%%%%%%%%%%%%%%%%%%%%%%%%%%%%%%%%%%%

See master bibliography (\texttt{bcm\_master.bib}) for complete citations.

\nocite{bcm_master}

